Under the alternative workflow all three prespecified contrasts excluded zero at
98.333\%, with the ordering independent $>$ clustered $>$ imperfect. Prevented
fractions were 76.5\% [74.1, 78.7], 53.7\% [49.7, 57.5] and 27.3\% [22.2, 31.9],
preventing 3.64, 2.31 and 1.17 acquisitions per 56-day run; bootstrap draws had
no undefined statistic. These results are conditional on the once-selected $p$.
Only the independent profile's baseline 95\% interval lay inside the descriptive
band; the clustered and imperfect holdout baselines were 2.62 and 2.63, below
2.7 (Table~\ref{tab:gates}), so the profiles were compared as fixed
configurations, not at a matched baseline. The clustered and imperfect
baselines lie within 0.005 of each other by chance in
these runs; they were not matched.

\emph{Baseline shortfalls (descriptive).} The 128-seed grid
estimates at m1 are imprecise and acquisitions are overdispersed (variance/mean
about 1.9); a descriptive bootstrap gives 2.46 [2.18, 2.77] (clustered) and 2.83
[2.51, 3.15] (imperfect) at m1. For imperfect clearing the selected $p$ is only
1.7\% above m1, yet the holdout baseline (2.63) was below the grid estimate
(2.83), so the grid point was probably high by chance. For clustered misses
grid error and log--log curvature cannot be separated (slope 2.5 below m1 and
3.35 above it). The protocol anticipated both kinds of error and neither is
propagated into the policy intervals. Section~\ref{sec:s-bias} bounds the
implied bias in $F$.

